% !TeX root = ../papers/nc-p3.tex
\section{量子点を閉軌道に持つ全三表の正則特殊化による閉性}
\label{sec:quantum-regular-closure}
\begin{theorem}\label{thm:quantum-regular-closure}
定理\ref{thm:quantum-universal-chambers}の三表の和、すなわち閉包に量子点を持つ窓全体の
射影閉包を、標準Koszul左右AS正則窓に制限すると、元の三表の和に一致する。
より詳しく、離散付値環上で幾何学的一般ファイバーの閉軌道が量子窓で、
特殊ファイバーが標準正則なら、一般の量子三角積は全て単元である。
特殊窓は半安定で、その閉軌道は三角積を還元して得た量子類である。
従って、固定した一つの量子類への全流入も標準正則軌跡で閉じている。
\end{theorem}
\begin{proof}[一般ファイバーの濾過と全順序での行列式]
有限拡大後、一般ファイバーの四位置の基底と一径数群を選び、最下重み部分を量子体積とできる。
実際、三表の構成では閉軌道の共通対角トーラスの余指標で基点に流入する。
各点で必要な有限個の基底係数は代数閉包上で選んだ後に有限拡大上で定義できる。
全行積が1の量子代表を使い、基点を各始点で同じ四文字順列の台を持つ体積とする。
余指標の四文字重みを $l_j$、体積の重みを $L=\sum_j l_j$ と書く。

各長さでは、小さい余指標重みを先頭とし、同じ重みでは任意に選んだ文字順序の左辞書順を使う。
元の関係空間の最下重み付きの随伴次数は量子関係空間である。実際、テンソルを一径数群で
量子体積に特殊化する間に両切断の階数6が保たれ、像六平面もその像に特殊化する。
従って六つの降順対に対し、量子二項関係にこれより高い重みの項を足した関係が取れる。
書換えは同じ重みでは転倒数を減らし、補正項では重みを真に上げる。長さごとに停止する。
昇順語数と標準次元の一致から全順序でこれがPBW基底となる。

補正項を使う回数は一つの書換え経路で $n(\max l_j-\min l_j)$ 以下である。
連続する二文字を他の二文字に替える一回の補正で転倒数の増加は $2n$ 以下であり、
同じ重みの書換えは転倒数を1減らす。従って全経路の書換え回数は一様に $O(n^2)$ である。
有限個の規則の係数から各経路の係数付値が $-Cn^2$ 以上となり、和でも同じ下界が保たれる。
定理\ref{thm:pbw-valuative-boundaries}の証明はこの重み優先辞書順にもそのまま適用できる。
実際、同じ両側の語を付加しても順序が保たれるため、交叉の先頭語は厳密降順であり、
経路と交叉の語の濾過で得る単項式完全列も同じである。

$w_i$ を降順四文字上でモニックにする倍率は背景と同じである。
文字 $j$ を除いた降順三文字が $W_{i+1,3}$ の座標を与えるので、
$1+3$ 対合の行 $a$、列 $j$ の成分の余指標重みは $L+l_a-l_j$ である。
$l_a<l_j$ の成分は零、等しい重みのブロックは量子体積の成分そのものである。
ゆえに行列式は背景の対角項の積に等しい。これは重みが重複しても変わらない。
\path{computations/verify_filtered_quantum_valuations.py} は三表の各重みと重複重みで、
全24順序の行列式と隣接交換による付値の差を検算する。
よって一般の変形係数によらず、前節の量子窓と全く同じ
$4\sum_{j<k}\operatorname{ord}c_{jk}\le0$ が全順序で必要である。
逆順と隣接交換により全ての $c_{jk}$ は単元である。
\end{proof}
\begin{proof}[量子閉軌道の還元と全被覆]
全交換係数が単元になった背景の各体積係数も、降順四文字でモニックにすれば単元である。
残る有限個のテンソル係数は全て重み $L$ を超えるので、余指標に一様化元の十分大きな正冪を
代入してから共通の重み $L$ の倍率を除くと、それらは全て極大イデアルに入る。
これで\emph{同じ一般窓}のもう一つの整模型が得られ、その特殊窓は剰余量子窓で標準正則である。
定理\ref{thm:regular-orbit-valuations}で元の正則模型と比較する。
四つの行列式の四乗根を取って特殊線形比較にすると、二つの原始最上元の倍率は単元である。
従って特殊窓での不変斉次式は共通の射影倍率を除いて同値であり、
元の特殊窓も半安定で剰余量子窓と同じ閉軌道を持つ。三角積の還元に対する同型は
定理\ref{thm:quantum-indexed-all}の文字置換を許す判定に従う。

最後にこの軌跡は構成可能である。量子体積とその基底移動は有限型データであり、
量子に流入する条件は三表それぞれで重みの下の座標が零、同じ重みの座標が背景に等しい
という条件と基底変換の存在で表せる。逆にこの条件で二切断階数6を満たせば
定理\ref{thm:quantum-universal-chambers}の範囲に入る。
その閉包の任意の正則点には、稠密な構成可能部分から曲線を取り正規化して上の付値判定を適用できる。
従ってその点も三表に入る。量子三角積の組を有限置換まで固定した場合も、
同じ構成可能性と還元の一意性から最後の主張を得る。
\end{proof}
